% ==============================================================================
% RESUMEN (español) — Página preliminar (APA 7 Nivel 1)
% Fuente: texto_formatter/resumen.txt (formateado en APA 7)
% ==============================================================================

\section*{Resumen}
\addcontentsline{toc}{section}{Resumen}

El presente proyecto tiene por objetivo desarrollar un sistema de información integral para optimizar el proceso de solicitud, recepción y selección de postulantes a la \enquote{Beca Comedor} en la Universidad Autónoma Gabriel René Moreno (UAGRM). Actualmente, la gestión de este beneficio presenta deficiencias debido a que la evaluación de solicitudes y perfiles se realiza de manera estrictamente manual. Esta situación genera cuellos de botella, tiempos de respuesta prolongados y un desgaste operativo significativo para el personal de la Dirección Universitaria de Bienestar Social y Salud (DUBSS). Adicionalmente, la población estudiantil carece de un canal digital unificado que facilite la postulación remota y la consulta transparente del menú diario.

Para solucionar esta problemática, se propone el desarrollo de una plataforma tecnológica estructurada en dos componentes: un portal web orientado a la gestión administrativa y una aplicación móvil para los estudiantes. El portal web permitirá a la DUBSS centralizar la información, gestionar perfiles y administrar el servicio de alimentación; por su parte, la aplicación móvil facilitará a los universitarios la solicitud de la beca y la consulta actualizada del menú del comedor.

El núcleo de la solución radica en la integración de Inteligencia Artificial (IA) para la preselección automatizada. Este modelo procesará los datos de los postulantes, descartando eficientemente aquellas solicitudes que no cumplan con los requisitos mínimos, reduciendo de forma masiva la carga de revisión humana. Complementariamente, se incorporará un módulo de Business Intelligence (BI) que proveerá tableros interactivos a las autoridades. Estas herramientas analíticas permitirán monitorear indicadores clave sobre la eficiencia de la preselección y el uso del comedor, facilitando una toma de decisiones informada. En conjunto, el proyecto moderniza y agiliza la administración del bienestar estudiantil.

\indent\textit{Palabras clave:} Sistema de Información, Comedor Universitario, Inteligencia Artificial, Business Intelligence, Automatización de Procesos, Aplicación Móvil, Gestión Administrativa.